% The spoken body of the homily for the Twenty-eighth Sunday in Ordinary Time,
% Year A. Continuous preaching, ready to read aloud: no heading, no direction
% to the preacher, no citation inside the speech. The loci, the relation to
% the reviewed interpretations, the audience, the word count and the pace
% estimate stand in the terminal note, which no one reads out. The blank lines
% are the paragraph breaks a speaker recovers his place by, and the six
% vertical spaces divide the speech into its seven movements: the invitation
% refused for a farm and a business; what the feast is that they refused; the
% guests from the roads; the apostle who is full and hungry; the wedding
% garment; this Mass as the feast now spread; and the one thing to do.
%
% Scripture is spoken in the Douay--Rheims, the leaf's study translation, and
% the speech names that Bible once, just before the first verse is quoted. No
% oration or antiphon is quoted in any Missal wording: the Collect, the Prayer
% over the Offerings and the Prayer after Communion are described, and of the
% two Communion antiphons only the first is named, as one of the two the
% Missal gives. The guest without the wedding garment is introduced as what
% Matthew's parable goes on to tell, so the speech holds whichever form of the
% Gospel was proclaimed.

A king makes a wedding for his son. The invitations have gone out, and when
the day comes the king sends his servants again. In the words of the old Douay
Bible: ``Tell them that were invited, Behold, I have prepared my dinner: my
beeves and fatlings are killed, and all things are ready. Come ye to the
marriage.''

And they do not come.

Saint John Chrysostom, preaching on this Gospel, puts the obvious question.
``For who would not choose to come to a marriage,'' he asks, ``and that a
King's marriage, and of a King making a marriage for a Son?'' Who turns down a
royal wedding?

Matthew tells us who. ``But they neglected and went their ways, one to his
farm and another to his merchandise.'' One to his farm. Another to his
business. Matthew does tell us of others, ``the rest'', who seized the king's
servants and killed them. But the first refusal is not violence. They simply
have somewhere else to be.

And there is nothing wrong with a farm. There is nothing wrong with a
business. Most of us have something like them: a job, a household, bills to
pay, work that has to be finished by Monday. These are not sins. They are the
ordinary goods of an ordinary life, and many of them are duties.

That is exactly what makes the parable so uncomfortable, and Chrysostom saw
it. ``The excuses seem to be reasonable,'' he says; ``but hence we learn,
though the things which hinder us be necessary, to set the things spiritual
at a higher price than all.'' And then, more sharply: ``not for press of
business, but from `making light of it,' they did not come.'' Matthew's word, in
the old translation, is \emph{neglected}. They did not hate the king. They simply did not
think his feast mattered as much as their own affairs. Saint Thomas Aquinas
says the same in a single sentence: outwardly they seemed to have a just
cause, but the Lord does not accept it, because nothing of this passing world
ought to keep anyone from coming to God.

So the question this Gospel puts to us is a quiet one. What is my farm? What
is my merchandise? What is the good, reasonable, necessary thing that I keep
allowing to stand between me and the King's invitation?

\bigskip

But first: what is being refused? What is this feast?

The first reading has already told us. ``The Lord of hosts shall make unto all
people in this mountain, a feast of fat things, a feast of wine.'' And on that
mountain, Isaiah says, ``He shall cast death down headlong for ever: and the
Lord God shall wipe away tears from every face.''

The Gospel tells us whose feast it is. Saint Jerome, explaining this parable,
says it plainly: this king is God almighty, and the wedding he makes is for
our Lord Jesus Christ and for the Church. Saint Gregory the Great asked what
that wedding is, and he answered with care.
The Father made a wedding for his Son, he says, when through the mystery of
the Incarnation he joined the holy Church to him. The feast is God's own Son,
made man, joined to his people.

And Gregory saw what the farm and the merchandise really cost. To go off to
the farm, he says, is to throw oneself into earthly labour beyond all measure;
to go off to the merchandise is to gape after the profits of worldly business.
A man who lives like that will not stop to weigh the mystery of the Lord's
Incarnation, or to live by it. That is the real refusal. Not a dinner missed,
but God made man passed over, because there was something else to do.

\bigskip

Now look at who does come. The king sends his servants into the highways, or,
as the Latin says more exactly, to the places where the roads go out. ``As
many as you shall find, call to the marriage.'' And the servants gather ``all
that they found, both bad and good: and the marriage was filled with guests.''

Gregory says of those roads that people very often come easily to God when no
prosperity accompanies them in their earthly doings. The people where the roads
went out had no farm to hurry back to. Nothing held them. And so they came.

There is hope here: no one is too far out on the road to be asked. And there
is a warning: our successes can be the very thing that keeps us at home.

\bigskip

Is the answer, then, to own nothing? Saint Paul gives a better answer in the
second reading. He writes to the Philippians, ``I know both how to be brought
low, and I know how to abound \ldots\ both to be full and to be hungry: both
to abound and to suffer need. I can do all things in him who strengtheneth
me.''

Paul is not without goods: the Philippians have sent him help. But nothing he
has holds him. Chrysostom says that plenty needs as much virtue as want does:
``For as want inclines us to do many evil things, so too doth plenty.''

Here is the guest who would have come. Not the man who has no farm, but the
man whose farm does not keep him.

And notice what the Philippians did with their goods. ``You have done well,''
Paul tells them, ``in communicating to my tribulation'': that is, in sharing
my trouble. The very things that kept the invited away became, in their hands,
a way of coming closer.

\bigskip

Matthew's parable goes on. The king comes in to see his guests, and finds one
man without a wedding garment. ``Friend,'' he says, ``how camest thou in
hither not having on a wedding garment?''

What is that garment? The Fathers give more than one answer. Saint Augustine
and Saint Gregory say: charity. Not baptism, Augustine says. No one comes to
God without it, but the good and the bad alike have been baptized. It is love. And Augustine tells his
people where to find it. ``Do not say; `we are too poor to have that garment.'
Clothe others, and ye are clothed yourselves. It is winter, clothe the naked.
Christ is naked; and He will give you that `wedding garment'.''

So the farm and the merchandise are not to be thrown away. They can be turned
into the wedding garment. What we have can keep us from the feast, or it can
clothe someone else, and so clothe us.

\bigskip

Jesus first told this parable in the temple, to the chief priests and the
elders of the people. But Gregory and Chrysostom, preaching it, turned it on
their own congregations, on the baptized, and so must we. Gregory told his
people that it is we who recline at the wedding of the Word. This Mass is the
King's feast, spread now.

At its beginning we prayed that God's grace would go before us and follow us,
and the Latin of that prayer asks that grace keep us intent on good works.
Gregory used that same word for the man who went off to his farm: he was
intent on earthly labour. We have asked to be intent on something else.

In a moment we will bring our gifts and ask that through this worship we may
pass over to heavenly glory. Then we will be fed with the Body and Blood of
Christ, and ask to be made partakers of the divine nature. One of the two
antiphons the Missal gives for Communion today says that the rich have known
want and hunger, but those who seek the Lord lack no good thing. Saint
Augustine said what the rich go without: ``Those rich then lack, they lack,
and what is heavier, they lack bread''; the bread of the one who said, ``I am
the Living Bread which came down from Heaven.''

\bigskip

So here is something to do this week. Name your farm. Name your merchandise.
Name the one good and reasonable thing that most often stands between you and
God: the work that spills into Sunday, the screen that swallows the hour you
meant for prayer, the worry about money that crowds out everything else. Do
not despise it. But take something from it this week and give it away: an
hour to someone who is alone, money to someone in need, a meal to someone who
is hungry. Clothe the naked, and you will be clothed.

And then come back. The servants are still being sent: to the farm, to the
shop, to the places where the roads go out. And the message has not changed.
``Behold, I have prepared my dinner \ldots\ and all things are ready. Come ye
to the marriage.''
